\documentclass{ranfr}

\title{ranfr.cls模板使用示例}
\author{RanFR}
\date{\today}

\begin{document}
\maketitle

\begin{abstract}
    在这个示例教程中，我们将演示如何使用ranfr.cls模板来创建一个简单的中文文档。这个文档包括数学公式、表格、图片、列表、代码和参考文献等内容。
\end{abstract}

\section{数学公式示例}

这是几个数学公式示例：

\begin{equation}
    E = mc^2
\end{equation}

\begin{equation}
    \int_0^1 x^2 \, dx = \frac{1}{3}
\end{equation}

\section{表格示例}

这是一个表格示例：

\begin{table}[h]
    \centering
    \caption{一个表格}
    \begin{tabularx}{\textwidth}{
            | >{\raggedright\arraybackslash}X
            | >{\centering\arraybackslash}X
            | >{\raggedleft\arraybackslash}X
            |}
        \hline
        项目（左） & 数量（中） & 单价（右） \\
        \hline
        A          & 5          & 20         \\
        \hline
        B          & 3          & 25         \\
        \hline
    \end{tabularx}
\end{table}

\section{列表示例}
\subsection{有序列表}
这是一个有序列表示例：
\begin{enumerate}
    \item 列表项 1
    \item 列表项 2
    \item 列表项 3
          \begin{enumerate}
              \item 子列表项 1
              \item 子列表项 2
          \end{enumerate}
\end{enumerate}

\subsection{无序列表}

这是一个无序列表示例：

\begin{itemize}
    \item 列表项 1
    \item 列表项 2
    \item 列表项 3
          \begin{itemize}
              \item 子列表项 1
              \item 子列表项 2
          \end{itemize}
\end{itemize}

\section{图片示例}

这是一张图片示例：

\begin{figure}[htbp]
    \centering
    \includegraphics[width=0.6\textwidth]{pictures/example.png}
    \caption{一张图片}
\end{figure}

\section{代码示例}

这是一个Python代码示例：

\begin{minted}{python}
    def hello_world():
        print("Hello, World!")

    hello_world()
\end{minted}

\section{伪代码示例}

这是一个简单的算法伪代码示例：

\begin{algorithm}[H]
    \caption{计算斐波那契数列的第n项}
    \KwIn{正整数$n$}
    \KwOut{第$n$项斐波那契数列的值}
    \SetKwProg{Fn}{Function}{\string:}{end}
    \Fn{Fibonacci($n$)}{
        \If{$n=1$ or $n=2$}{
            \Return 1;
        }
        $F_1\leftarrow 1$; $F_2\leftarrow 1$;
        \For{$i\leftarrow 3$ \KwTo $n$}{
            $F_i\leftarrow F_{i-1} + F_{i-2}$;
        }
        \Return $F_n$;
    }
\end{algorithm}

\end{document}
